\documentclass[sigconf,9pt]{acmart}
% Candidate intentionally has no author block: verified owner metadata is absent.
% This is a nonanonymous-track preparation candidate, not an upload-ready PDF.
\setcopyright{none}
\settopmatter{printacmref=false}
\acmConference[SIGMOD 2027 Demonstrations]{SIGMOD 2027 Demonstrations}{}{}
\acmYear{2027}
\acmDOI{}
\acmISBN{}
\title{CDC Replay Correctness Lab: Making Stable but Incorrect State Visible}
\begin{abstract}
Change-data capture consumers can produce a stable but incorrect materialization:
identity deduplication does not order distinct changes, and finite identity memory
may hide conflicting historical payloads without changing the latest state.
CDC Replay Correctness Lab is a local, interactive demonstration of these
 distinctions. Reviewers select ordering, replay policy, checkpoint and interruption
boundaries and ID retention, inspect event-by-event decisions, and compare actual
state with an independent canonical oracle. A guaranteed stale-delete preset
reproduces baseline failure without searching seeds. Independent execution of the
original artifact reproduced its 126 principal and 36 bounded observations. New
isolated controls show that all three policies pass duplicate-only and ordered
delete/reinsert workloads; an ID-only policy still fails the stale-delete case.
A separate probe-order study exposes why retention-dependent conflict detection
cannot be inferred from final-state equality. The contribution is an inspectable,
reproducible educational and diagnostic artifact, not a new CDC algorithm or an
exactly-once guarantee.
\end{abstract}
\ccsdesc[500]{Software and its engineering~Software testing and debugging}
\ccsdesc[300]{Information systems~Data replication tools}
\keywords{change-data capture, replay, correctness oracle, interactive demonstration}
\begin{document}
\maketitle
\section{Problem and demonstration contribution}
A downstream materializer can receive a historical deletion after a newer reinsertion.
Keeping only unique event IDs does not prevent this stale overwrite when the IDs
are distinct. Moreover, a second replay may leave an already incorrect state unchanged.
A demonstration that reports only duplicate counts or redelivery stability can
therefore conceal the failure it is supposed to explain.

The lab turns these counterexamples into an audience-controlled comparison of three
simple policies. It separates \emph{final-state equality}, \emph{duplicate recognition},
\emph{conflict detection}, and \emph{recovery attempts}. A reviewer can switch from
whole-input validation to online processing and expire identities while retaining
versioned tombstones. This reveals both containment and missing detection rather
than displaying only favorable outcomes. The artifact makes established semantics
observable; it introduces no new fault-tolerance algorithm and reports no user-study
or production result.

The event contract assumes immutable identity, an authoritative comparable per-key
integer version, a string payload and a deletion flag. Higher versions determine the
intended state; a latest deleted record remains a tombstone until superseded. Real sources require source-specific
ordering contracts: Debezium's PostgreSQL documentation describes LSN-bearing change
events and offset-based restart~\cite{debezium}. We neither derive comparable versions
from arbitrary CDC records nor implement that connector.

\section{Architecture and correctness model}
\begin{figure}[t]
\centering
\fbox{\parbox{0.93\columnwidth}{\centering
Canonical synthetic history $\longrightarrow$ independent batch oracle\\[3pt]
$\downarrow$ scenario / ordering / failure construction\\[3pt]
Arrival $\mid$ ID-only $\mid$ version-aware processors\\[3pt]
$\downarrow$ logical checkpoint, restore, redelivery\\[3pt]
Per-event trace $+$ counters $+$ actual/oracle comparison}}
\caption{The canonical specification and replay transitions are separate. The interface
shares the tested model with the V2 experiment runner.}
\Description{Dataflow separates the canonical batch oracle from scenario generation and three replay policies, then compares their traces and final states.}
\label{fig:architecture}
\end{figure}
The system has a Python HTTP server from the standard library, embedded HTML and JavaScript,
a replay model, two research runners and artifact validators. It needs no external
service or front-end download. The installed wheel serves the same interface from
outside the checkout. Requests describe a bounded synthetic scenario; they do not
upload arbitrary client data. State is local to one request.

The independently implemented oracle groups the canonical history by key and selects the maximum
version. Its identity and key/version ambiguity checks are separately implemented;
it calls no materializer transitions or shared processor conflict validator. It is
computed before injected delivery is processed. Hand-authored tests assert stale
updates, deletion/reinsertion and ambiguous identities/versions. This is a test
specification, not an external database authority or formal proof.

Arrival-order replay overwrites on each event. ID-only replay suppresses remembered
IDs but otherwise overwrites. Version-aware replay also refuses lower-version state
writes and conflicting equal-version values. Principal ID-based policies use an unlimited ID ledger; arrival replay keeps none. A shared preflight validator rejects conflicting input before any of the
three policies executes; that rejection must not be called online conflict detection.

At checkpoint offset $c$, logical state and ID memory are JSON-serialized. Execution
continues to interruption $i$, discards later attempts, restores the earlier snapshot
and replays from $c$. These are simulated logical checkpoints: no disk write,
fsync, transactional source/sink coupling or real crash is provided. Distributed
consistent cuts are a substantially stronger problem~\cite{snapshots}.

For completed recovery runs, the trace marks events attempted before interruption
that are discarded when the snapshot is restored. Monotonicity counters
separate attempted writes, committed-prefix writes and recovered writes. Restore can
lower a visible version because an earlier snapshot is reinstated; it is not counted
as a stale materialization write. A completed run's correctness requires whole-state
oracle equality, independently of its redelivery stability.

\section{Live demonstration experience}
\begin{figure}[t]
\centering
\fbox{\parbox{0.91\columnwidth}{
\textbf{Stale-delete preset: one key, native order}\\[3pt]
\texttt{b0:v4}: live v4 $\longrightarrow$ \texttt{b0:v3}: delete v3\\[4pt]
\begin{tabular}{lll}
Policy & Actual & Oracle\\\hline
Arrival & deleted v3 & live v4\\
ID-only & deleted v3 & live v4\\
Version-aware & live v4 & live v4
\end{tabular}\\[4pt]
\textbf{Inspect:} strategy switch; wrong-key reason;
trace position; before/after state; stale-write decision.
}}
\caption{Readable typeset view of verified preset outputs, not a browser screenshot.
The working interface exposes these states and interactions; the artifact retains
the full Chromium screenshot. Distinct IDs do not protect against stale deletion.}
\Description{Two deliveries, live version 4 then deleted version 3, leave arrival and ID-only replay wrong while version-aware replay matches the oracle.}
\label{fig:ui}
\end{figure}
The interface exposes nine scenarios, policy selection, native, reversed, seeded or version-based
ordering, logical checkpoint and interruption offsets, and finite or unlimited ID
capacity. Side-by-side states and wrong-key explanations identify a missing key,
stale version, unexpected version, or payload/tombstone mismatch. Trace rows show
phase, source offset, event identity, applied, duplicate, stale or conflict decisions and
ledger size. A slider reveals state after the selected transition. Invalid controls
produce an explicit error; rejected or conflict-stopped runs have no completed
correctness verdict.

The five-minute walkthrough follows five interactions:
\begin{enumerate}
\item Click the stale-delete preset: one key receives live v4 under one ID and deleted
v3 under another. Arrival-order replay visibly fails within the automated one-minute
check. This case is deterministic and requires no seed search.
\item Select ID-only, then version-aware. ID-only remains wrong; version-aware keeps
live v4. Inspect the stale decision and before/after state with the trace controls.
\item Choose recovery/redelivery, three keys, checkpoint 4 and interruption 8.
Inspect four discarded attempts and four replayed lost attempts; compare final state
and the distinct redelivery-stability indicator.
\item Choose ledger expiry, 16 keys, online validation and capacity 4. Latest state
matches the oracle, yet an injected conflicting stale ID payload is not detected;
the interface explicitly warns about this distinction.
\item Clear the capacity for unlimited retention. The same conflict is detected and
processing stops. Switch to whole-input validation to see rejection before replay.
\end{enumerate}
The allocation is a presenter script, not an observed five-minute usability study.
HTTP integration tests cover preflight rejection. Real Chromium automation executes
the preset, policy switches, trace stepping, recovery and finite and unlimited ledger
interactions with no JavaScript errors. The browser uses the built wheel outside the repository.

\section{Evidence and isolated controls}
The protocol was repository-committed before implementation and measurements.
A new execution of the unchanged original runner reproduced all non-timing fields
of its 126 principal and 36 bounded observations. Artifact validation checked original
and new SHA-256 manifests, CSV/JSON field parity, source and configuration hashes and
row counts. This is independent execution, not replication by a separately authored
implementation. Timing samples were intentionally not required to match.

Original duplicate and delete/reinsert conditions also shuffled input; they were
compound workloads. V2 adds truly isolated controls and records a new 144-row matrix:
8/64 keys, three seeds, eight conditions, three policies, five fresh timed executions.
Some ordered conditions do not vary with seed; rows are not independent statistical
samples. Conflicting-ID and same-version conditions yield 12 rejected inputs per
policy. The remaining 36 completed observations per policy are summarized below.

\begin{table}[t]
\caption{V2 outcomes per policy: 36 completed and 12 preflight-rejected inputs.}
\label{tab:results}
\centering
\begin{tabular}{lrrr}\toprule
Policy & Correct & Incorrect & Rejected\\\midrule
Arrival & 24 & 12 & 12\\
ID-only & 30 & 6 & 12\\
Version-aware & 36 & 0 & 12\\\bottomrule
\end{tabular}
\end{table}
All policies pass all six configurations each of ordered live input, duplicate-only,
ordered delete/reinsert and ordered checkpoint recovery. Thus neither duplicates
alone nor an ordered deletion/reinsertion sequence establishes an advantage for
version guards in this final-state model. Arrival and ID-only fail all six stale-delete
configurations. Reverse post-recovery redelivery additionally makes arrival replay
wrong; remembered identities contain that case for ID-only.

A separate 72-row online matrix crosses policy, capacities of 0, 4, 16 or unlimited, three
recorded seeds and two probe orders. State and tombstones remain unbounded. In
\emph{conflict-first}, version-aware replay detects the expired-ID conflict only with
unlimited memory. In \emph{exact-then-conflict}, capacities 4/16 detect it after exact
stale redelivery reinserts the ID; capacity 0 still misses it. Each configuration is
repeated over three seeds, but its ordered bootstrap is identical. A version guard
can preserve latest state while losing historical-payload consistency checks.

Five timing samples and per-row medians are retained for every observation. Principal
run timing includes validation, snapshot copies, serialization and event traces;
expiry timing covers the materializer bootstrap and probes, without traces or
checkpointing. Both exclude data construction, oracle/scoring and output I/O. Trace copying changes measured work, so
V2 timings are not used as speedups over the original runner. One host, small synthetic
streams and early rejection preclude throughput, cloud-cost or production conclusions.

\section{Related work, limits and reproducibility}
MillWheel studies persistent state and fault-tolerant streaming~\cite{millwheel}; the
Dataflow Model treats correctness, latency and cost for unordered input~\cite{dataflow}.
This lab does not reproduce either system. Its reference-specification approach sits
within established model-based testing~\cite{mbt}; it does not benchmark against that taxonomy or introduce a new test-oracle method. Interactive
algorithm visualizations such as RaftScope~\cite{raftscope} predate this work. Our
specific demonstration joins CDC stale-delete counterexamples, attempted and recovered
trace evidence and probe-order-dependent ID expiry, rather than claiming invention
of visualization, idempotency, versions or checkpoints.

The model has no distributed scheduler, external CDC connector, durable sink,
concurrency, source epochs or tombstone garbage collection. Stable IDs and authoritative
versions are assumptions. Equal versions cannot resolve ambiguous data. Final-state
equality does not measure external side-effect multiplicity or imply exactly-once
processing. A finite ID ledger is not a bound on total state. Claims are limited to
retained deterministic workloads; independent implementation replication and a
broader prior-art review remain outstanding.

The artifact provides commands for unit tests, frozen V4/V5 validation, and the V2
runner with fresh output paths. Start the interface using Python's \texttt{-m}
option with module \texttt{pipeline\_reliability.}\allowbreak\texttt{cdc\_interactive}.
Retained artifacts record full source commit, dirty status, Python and platform, seeds,
configurations, all samples, and hashes of the workload, source and protocol. The retained versioned evidence, a five-minute walkthrough, and fresh-output commands
are available in the public artifact~\cite{artifact}. Logical checkpoint behavior
is not durable crash recovery; preflight rejection is not online detection.
\bibliographystyle{ACM-Reference-Format}
\bibliography{references}
\end{document}
